% SGA 3, Expose VI_A, Loop 1 English reconstruction.
% Scope: section 5 heading and paragraph 5.1 in full;
% paragraph 5.2 is excluded.
% Authority: Polo--Gille Exp6A-13oct24.pdf, component-local pages 30--31
% (combined-reader pages 333--334; printed pages 322--323).
% OCR and the Reinhold English Markdown are comparison witnesses only.
% Loop 2 replaces the exact 180-dpi source crop with a native tikz-cd
% reconstruction controlled against local p. 30,
% x=520,y=1142,w=530,h=305 on a 1488x2105 render.

\section*{5. 与第 IV 篇的联系及其推论}
\phantomsection\label{sec:sga3-VIA-5}
\addcontentsline{toc}{section}{5. 与第 IV 篇的联系及其推论}

\medskip
\phantomsection\label{par:sga3-VIA-5.1}%
\noindent\textbf{5.1.---}%
\footnote{\label{note:sga3-VIA-49}编者注: 本节标题已作修改,
原文标题为“补充”.}
沿用 \S~3 的记号和\hyperref[thm:sga3-VIA-3.2]{定理~3.2} 的假设.
于是有如下交换图:
\begin{center}
  \phantomsection\label{diag:sga3-VIA-5.1-right-action}%
  \begin{tikzcd}[row sep=huge,column sep=huge]
    F\times_A G\times_A G
      \arrow[r,"F\times\nu"]
      \arrow[d,shift right=0.9ex,"\operatorname{pr}_2\times G"']
      \arrow[d,shift left=0.9ex,"\lambda\times G"] &
    F\times_A G
      \arrow[d,shift right=0.9ex,"\operatorname{pr}_2"']
      \arrow[d,shift left=0.9ex,"\lambda"] \\
    G\times_A G
      \arrow[r,"\nu"]
      \arrow[d,"p\times G"'] &
    G
      \arrow[d,"p"] \\
    (F\backslash G)\times_A G
      \arrow[r,dashed,"\rho"] &
    F\backslash G
  \end{tikzcd}
\end{center}
其中满足
\[
  \operatorname{pr}_2\circ(F\times\nu)
    =
  \nu\circ(\operatorname{pr}_2\times G),
  \qquad
  \lambda\circ(F\times\nu)
    =
  \nu\circ(\lambda\times G).
\]
此外, 由于 \(G\) 在 \(A\) 上平坦,
由 \hyperref[par:sga3-VIA-4.7]{4.7} 可知左边的竖直序列正合.
因此 \(\nu\) 诱导一个 \(A\)-概形态射
\[
  \rho:(F\backslash G)\times_A G\longrightarrow F\backslash G.
\]

立即可验, 态射 \(\rho\) 使 \(G\) 在 \(F\backslash G\) 上右作用;
此外, 典范态射 \(G\to F\backslash G\) 与 \(G\) 在 \(G\)
及 \(F\backslash G\) 上的右作用可交换.%
\footnote{\label{note:sga3-VIA-50}编者注: 下文为增补内容.}
这证明了\hyperref[concl:sga3-VIA-3.2-ii-prime]{定理~3.2}(ii\(^{\prime}\)) 的第一个断言.
因此由 \SGARefLink{prop:sga3-VIA-2.5.4}{2.5.4},
\[
  X=F\backslash G
\]
的各个连通分支均为有限型、不可约, 且维数相同.

为计算这个维数, 可以假设 \(A=k\), 且 \(k\) 代数闭.
由\SGARefLink{sga3:i:scholium:2-3-3-1}{第~I~篇 2.3.3.1},
\(k\)-点 \(p(e)\) 的稳定子由纤维
\[
  H=p^{-1}\bigl(p(e)\bigr).
\]
代表. 由于 \(F\backslash G\) 是赋环空间范畴中 \(G\) 对 \(F\) 的商,
该纤维的底空间为 \(u(F)\). 由于 \(\ker(u)\) 有限, 得到
\[
  \dim H=\dim u(F)=\dim F.
\]
由 \hyperref[concl:sga3-VIA-2.5.4-ii]{2.5.4(ii)} 可得
\[
  \dim X=\dim G-\dim F.
\]
这证明了\hyperref[concl:sga3-VIA-3.2-ii-prime]{定理~3.2}(ii\(^{\prime}\)),
因而也证明了\hyperref[concl:sga3-VIA-3.3.2-iv]{定理~3.3.2(iv)}.
